\chapter{ANALISIS MASALAH}


Tuliskan di sini \textbf{Pedoman Penulisan Analisis Masalah:} \\
Lakukan identifikasi dan analisis mendalam terhadap masalah-masalah spesifik yang melatarbelakangi kebutuhan akan perangkat lunak inovatif ini. Uraikan masalah tersebut menjadi poin-poin yang terstruktur agar mudah dipahami.\\
Sajikan data pendukung yang valid (misalnya hasil survei awal, observasi langsung, wawancara dengan calon pengguna, atau referensi studi kasus sejenis) untuk membuktikan bahwa masalah tersebut benar-benar ada dan krusial.\\
Jelaskan dampak buruk atau konsekuensi negatif yang akan terjadi di masa mendatang bagi pengguna atau perusahaan/instansi terkait jika permasalahan tersebut tidak segera diselesaikan atau dibiarkan tanpa solusi teknologi.

